\documentclass[12pt]{article}
\usepackage[utf8]{inputenc}
\usepackage[T1]{fontenc}
\usepackage[spanish]{babel}
\usepackage{geometry}
\usepackage{hyperref}
\usepackage{enumitem}
\usepackage{verbatim}
\geometry{margin=2.5cm}

\title{Instalaci\'on de Docker en Ubuntu 20.04 LTS}
\author{Gu\'ia basada en la configuraci\'on real de esta m\'aquina}
\date{25 de abril de 2026}

\begin{document}
\maketitle

\section*{1. Instalar dependencias base}
Se actualizaron los índices y se instalaron las utilidades necesarias para agregar el repositorio oficial de Docker.

\begin{verbatim}
sudo apt-get update
sudo apt-get install -y ca-certificates curl gnupg lsb-release
\end{verbatim}

Estas herramientas permiten descargar la clave GPG del repositorio oficial y registrar el origen de paquetes de Docker de forma segura.

\section*{2. Agregar la clave GPG de Docker}
Para que \texttt{apt} confíe en el repositorio oficial, se creó el directorio de llaves y se instaló la clave de firma de Docker.

\begin{verbatim}
sudo install -m 0755 -d /etc/apt/keyrings
curl -fsSL https://download.docker.com/linux/ubuntu/gpg \
  | sudo gpg --dearmor -o /etc/apt/keyrings/docker.gpg
sudo chmod a+r /etc/apt/keyrings/docker.gpg
\end{verbatim}

Durante este paso, la clave ya existía en la máquina y se sobrescribió para dejarla consistente con el repositorio oficial.

\section*{3. Registrar el repositorio oficial}
En lugar de depender del paquete \texttt{docker.io} de Ubuntu, se agregó el repositorio oficial de Docker para obtener las versiones actuales de Engine, Compose y Buildx.

{\small
\begin{verbatim}
repo="deb [arch=$(dpkg --print-architecture) \\
signed-by=/etc/apt/keyrings/docker.gpg]"
repo="$repo https://download.docker.com/linux/ubuntu $(lsb_release -cs) stable"
echo "$repo" | sudo tee /etc/apt/sources.list.d/docker.list > /dev/null
sudo apt-get update
\end{verbatim}
}

En esta máquina la distribución reportó \texttt{focal}, por lo que el repositorio quedó apuntando a Ubuntu 20.04.

\section*{4. Quitar paquetes anteriores de Ubuntu}
Antes de instalar Docker desde el repositorio oficial, se eliminaron los paquetes que venían del repositorio de Ubuntu para evitar conflictos.

\begin{verbatim}
sudo apt-get remove -y docker.io containerd runc
\end{verbatim}

Esto liberó la instalación previa y evitó mezclar binarios de distinta procedencia.

\section*{5. Instalar Docker completo}
Luego se instaló el conjunto oficial de paquetes:

\begin{itemize}[leftmargin=*]
  \item \texttt{docker-ce}: motor de Docker.
  \item \texttt{docker-ce-cli}: cliente de línea de comando.
  \item \texttt{containerd.io}: runtime de contenedores.
  \item \texttt{docker-buildx-plugin}: soporte para Buildx.
  \item \texttt{docker-compose-plugin}: soporte para \texttt{docker compose}.
  \item \texttt{docker-ce-rootless-extras}: utilidades para modo rootless.
\end{itemize}

\begin{verbatim}
sudo apt-get install -y docker-ce docker-ce-cli containerd.io \
  docker-buildx-plugin docker-compose-plugin docker-ce-rootless-extras
\end{verbatim}

En esta m\'aquina, \texttt{apt} instal\'o adem\'as \texttt{slirp4netns}, que es una dependencia usada por el modo rootless.

\section*{6. Habilitar el servicio}
Se activ\'o Docker para que arranque autom\'aticamente con el sistema.

\begin{verbatim}
sudo systemctl enable --now docker
\end{verbatim}

Con esto, el servicio qued\'o activo y habilitado al inicio.

\section*{7. Verificar la instalaci\'on}
Finalmente se comprob\'o que las herramientas respondieran correctamente y se ejecut\'o un contenedor de prueba.

\begin{verbatim}
docker --version
docker compose version
docker buildx version
systemctl is-active docker
systemctl is-enabled docker
docker run --rm hello-world
\end{verbatim}

La verificaci\'on fue exitosa: Docker report\'o la versi\'on \texttt{28.1.1}, Compose report\'o \texttt{v2.35.1}, Buildx report\'o \texttt{v0.23.0} y el contenedor \texttt{hello-world} confirm\'o que el motor funciona correctamente.

\section*{8. Resultado final}
La instalaci\'on correcta en esta computadora qued\'o con esta combinaci\'on:

\begin{itemize}[leftmargin=*]
  \item Ubuntu 20.04.6 LTS.
  \item Repositorio oficial de Docker para \texttt{focal}.
  \item Docker Engine 28.1.1.
  \item Docker Compose Plugin 2.35.1.
  \item Docker Buildx Plugin 0.23.0.
  \item Servicio \texttt{docker} activo y habilitado.
\end{itemize}


\end{document}
